\section{From conserved quantities to invariant packets}
\label{sec:packets}
\status{Implemented in \code{invariant_space.py}; exact identity checker and
recorded experiments included.}

\subsection{Preservation of a zero set is weaker than conservation}
Let a state be $x\in\Q^n$ and let $T_1,\ldots,T_r:\Q^n\to\Q^n$ be polynomial
transitions. Every transition is allowed at every state in the initial fragment.
Let the initial states be the union of the images of polynomial maps
$\sigma_1,\ldots,\sigma_s$. Parameters of these maps range over all rational
values. The implementation pads parameter maps with unused variables so all
polynomials use one common arity.

A conserved polynomial satisfies $p(T_a(x))=p(x)$ everywhere. An invariant
equality only requires that $p=0$ remain true on reachable states. Requiring
conservation is therefore unnecessarily strong. With several polynomials, a
sufficient preservation condition is
\begin{equation}
\label{eq:constant-packet}
 P(\sigma_b(u))=0,
 \qquad P(T_a(x))=M_a P(x),
\end{equation}
where $P=(p_1,\ldots,p_k)^{\mathsf T}$ and $M_a$ is a rational matrix. No
individual $p_i$ need be conserved, and $M_a$ need not be invertible, positive,
or stable in any analytic sense.

\begin{theorem}[Packet soundness]
\label{thm:packet}
If the polynomial identities in \eqref{eq:constant-packet} hold, then
$p_i(x)=0$ for every reachable state $x$ and every $i$.
\end{theorem}
\begin{proof}
The initial identities establish the claim at a state $\sigma_b(u)$. If
$P(x)=0$ and the next transition is $T_a$, then
$P(T_a(x))=M_a P(x)=M_a0=0$. Induction on the length of a finite execution gives
the result. No termination statement about an infinite execution is required.
\end{proof}

\begin{example}[A packet that conservation misses]
\label{ex:cycle}
Take
\[
 T(x,y,z)=(y,z,x),\qquad (x_0,y_0,z_0)=(t,t,t).
\]
With $p=x-y$ and $q=y-z$,
\[
 \begin{pmatrix}p\circ T\\q\circ T\end{pmatrix}
 =\begin{pmatrix}0&1\\-1&-1\end{pmatrix}
  \begin{pmatrix}p\\q\end{pmatrix}.
\]
Thus $x=y=z$ is preserved. In contrast, a conserved linear polynomial has equal
coefficients on $x,y,z$; requiring that it vanish at every $(t,t,t)$ forces
those coefficients and the constant term to be zero. The conserved-linear
subspace is zero, while the invariant-linear subspace has dimension two.
The recorded prototype returns the equivalent basis $(z-x,y-x)$, not the
hand-chosen basis displayed above.
\end{example}

This is a strict increase in what the \emph{selected template method} can
certify. It is not a demonstration that Lean's existing tactics cannot prove
this elementary example with suitable induction and hypotheses.

\subsection{Compute the entire admissible space}
Fix a finite-dimensional template space
\[
 V_d=\{p\in\Q[X_1,\ldots,X_n]:\degp p\le d\},
 \qquad N=\dim V_d=\binom{n+d}{d}.
\]
A sparse template can replace all monomials, but the full degree template has
a particularly clean completeness statement. Polynomial pullback
$T_a^*:p\mapsto p\circ T_a$ is a \emph{linear operator in $p$}, even when
$T_a$ is nonlinear. Its output may have a higher degree; that is not a problem
for exact coefficient comparison.

First compute
\[
 W_0=\{p\in V_d:p\circ\sigma_b=0\text{ as a polynomial for every }b\}.
\]
Substitution of each template monomial into each initial map gives a matrix of
rational coefficients. Its kernel is $W_0$. This is not interpolation over a
sample of initial states.

For a current space $W$, apply the restriction
\begin{equation}
\label{eq:constant-restriction}
 W\longleftarrow W\cap (T_a^*)^{-1}(W).
\end{equation}
Cycle through all transitions until a full pass makes no dimension decrease.
Because every update only removes vectors, equality of dimensions implies
equality of spaces; comparing pretty-printed bases is unnecessary.

\subsection{The actual linear system}
Let $b_1,\ldots,b_k$ be a basis of the current $W$. Form coefficient columns for
\[
 b_1\circ T_a,\ldots,b_k\circ T_a,-b_1,\ldots,-b_k.
\]
A null vector $(c,h)$ says
\[
 \sum_i c_i(b_i\circ T_a)=\sum_j h_j b_j.
\]
Project a basis of this kernel to its first $k$ coordinates and take the span
of the resulting polynomials $\sum_i c_i b_i$. This is exactly the new space in
\eqref{eq:constant-restriction}. Kernel vectors with $c=0$ contribute no
polynomial and disappear during row reduction.

\begin{lstlisting}
W := kernel(initial_substitution_constraints(V_d))
repeat:
    old_dimension := dimension(W)
    for T in transitions:
        A := coefficient_columns(pullback(T, basis(W)))
        B := coefficient_columns(basis(W))
        K := nullspace([A | -B])
        W := span(project_left(K) * basis(W))
until dimension(W) = old_dimension
solve coordinates of each pullback in the final basis
return basis(W), transition_matrices
\end{lstlisting}

The source uses rational row reduction, deterministic monomial order, and
canonical rational coefficients. It is deliberately small rather than an
optimized symbolic linear-algebra library.

\begin{theorem}[Greatest pullback-stable subspace]
\label{thm:greatest-constant}
The restriction algorithm terminates. Its result is the greatest subspace
$W\subseteq W_0$ for which $T_a^*W\subseteq W$ for every transition $a$.
There are at most $N$ strict dimension decreases.
\end{theorem}
\begin{proof}
Every update is an intersection with a linear preimage, hence is a subspace of
the previous space. Strict inclusion decreases dimension. A full pass with no
decrease leaves every restriction unchanged, establishing closure.
If $U\subseteq W_0$ is any common pullback-stable subspace, then $U\subseteq W$
before an update implies $T_a^*U\subseteq U\subseteq W$, and hence
$U\subseteq W\cap(T_a^*)^{-1}(W)$ afterwards. Induction over updates shows
$U$ is contained in the final space.
\end{proof}

\begin{theorem}[Exact degree-bounded equality invariants for affine transitions]
\label{thm:affine-complete}
If every transition is affine, the final space equals the set of all
polynomials of degree at most $d$ that vanish on the reachable states defined
above.
\end{theorem}
\begin{proof}
Let $E$ be that set. It is a subspace of $W_0$. For $p\in E$, an affine
substitution preserves the degree bound, and $p\circ T_a$ vanishes on every
reachable state because each successor is reachable. Thus $T_a^*E\subseteq E$,
so Theorem~\ref{thm:greatest-constant} gives $E\subseteq W$. Conversely the final
basis has a packet certificate, and Theorem~\ref{thm:packet} gives $W\subseteq E$.
\end{proof}

This is a lifted, degree-bounded linear-invariant computation, in the lineage of
Karr-style affine relation analysis. It is not the much more general algebraic
algorithm for all polynomial invariants of affine programs studied by
Hrushovski, Ouaknine, Pouly, and Worrell~\cite{karr,affine-invariants}.

\subsection{A nonlinear case the constant packet still finds}
Consider $T(n,s)=(n+1,s+n^2)$ with initial state $(0,0)$. In degree three the
successive recorded dimensions are
\[
 9\;\longrightarrow\;5\;\longrightarrow\;3\;\longrightarrow\;2
 \;\longrightarrow\;1\;\longrightarrow\;1.
\]
The surviving basis polynomial is
\[
 p(n,s)=s-\frac{n}{6}+\frac{n^2}{2}-\frac{n^3}{3}.
\]
Its checked pullback is $p\circ T=p$. Clearing denominators yields
$6s=n(n-1)(2n-1)$. Forge's existing recurrence worker also targets this kind of
example; it is included here to test compatibility with a known successful
fragment, not as a new capability claim.

\subsection{An explicit limitation, not a reason to abandon the method}
For $T(x,y)=(x^2,y^2)$ and initial states $(t,t)$, the true linear invariant is
$x-y=0$. But
\[
 (x-y)\circ T=(x+y)(x-y),
\]
which is not in the one-dimensional span of $x-y$. At $d=1$ the constant-packet
algorithm returns the zero space. The failure comes from its certificate
language, not from the absence of a linear invariant. The next section removes
this particular obstruction without a bilinear solver.

\section{Invariant ideals with bounded multipliers}
\label{sec:ideals}
\status{Implemented in \code{ideal_space.py}. Polynomial multiplier certificates
and negative controls are included.}

\subsection{A monotone relaxation of the packet condition}
For a nonnegative multiplier-degree bound $e$, define
\begin{equation}
\label{eq:bounded-ideal}
 J_e(W)=\Span_{\Q}\{m p : p\in W,\ m\text{ a monomial},\ \degp m\le e\}.
\end{equation}
This is a finite-dimensional \emph{truncated multiple space}, not necessarily
an ideal: it need not be closed under further multiplication. In terms of a
basis $p_1,\ldots,p_k$ of $W$, it consists precisely of sums
$\sum_j h_jp_j$ with $\degp h_j\le e$. The notation $J_e$ is convenient, but
using the full ideal's closure properties here would be a mistake.

Replace the restriction in the previous section by
\begin{equation}
\label{eq:ideal-restriction}
 F_{a,e}(W)=W\cap(T_a^*)^{-1}\bigl(J_e(W)\bigr).
\end{equation}
This operation is monotone: $U\subseteq W$ implies
$J_e(U)\subseteq J_e(W)$, and therefore $F_{a,e}(U)\subseteq F_{a,e}(W)$.
It is also deflationary, since $F_{a,e}(W)\subseteq W$.

The final certificate becomes
\begin{equation}
\label{eq:poly-packet}
 p_i\circ\sigma_b=0,
 \qquad p_i\circ T_a=\sum_{j=1}^k h_{aij}p_j,
 \qquad \degp h_{aij}\le e.
\end{equation}
The same induction proof as before establishes soundness, because every term
on the right vanishes when all $p_j$ vanish. Coefficients may be polynomials in
the current state. They do not have to be invariant or bounded in magnitude.

\subsection{Why this does not solve a bilinear system}
Forge correctly notes that writing unknown generators $p_i$ and unknown
multipliers $h_{ij}$ simultaneously produces bilinear coefficient equations.
The algorithm here does not solve that system directly. It starts with
\emph{all} initially valid polynomials in a fixed finite template. At every
restriction step the current generator basis is already known. Consequently,
all products of that basis with bounded monomials are known polynomials, and
finding which pullbacks belong to their span is ordinary linear algebra.

For current basis $b_1,\ldots,b_k$ and multiplier monomials
$m_1,\ldots,m_K$, form columns
\[
 b_i\circ T_a\quad(1\le i\le k),
 \qquad -m_\ell b_j\quad(1\le j\le k,\ 1\le\ell\le K).
\]
Take the kernel and project to the first $k$ coordinates as before. This
computes \eqref{eq:ideal-restriction} exactly. After convergence, solve for the
coordinates of each $b_i\circ T_a$ in the final product list and regroup the
coefficients into polynomials $h_{aij}$.

\begin{lstlisting}
W := all initially-zero polynomials in V_d
repeat:
    old_dimension := dimension(W)
    for T in transitions:
        products := [m * p for p in basis(W), deg(m) <= e]
        columns := pullback(T, basis(W)) ++ (-products)
        K := nullspace(coefficient_matrix(columns))
        W := span(left_projection(K) * basis(W))
until dimension(W) = old_dimension
H := coordinates(pullbacks(final_basis), final_products)
return final_basis, regroup_as_polynomial_matrices(H)
\end{lstlisting}

This is a change in the search problem, not a linearization that drops products
of unknown quantities and hopes they can be repaired later. Every certificate
identity is checked with all products present.

\begin{theorem}[Greatest bounded-multiplier generator space]
\label{thm:ideal-greatest}
The algorithm terminates and returns the greatest subspace $W\subseteq W_0$
such that $T_a^*W\subseteq J_e(W)$ for every $a$.
\end{theorem}
\begin{proof}
The sequence descends inside finite-dimensional $V_d$, so it stabilizes.
Stability of a complete transition pass is exactly the required inclusion for
each transition. Let $U\subseteq W_0$ satisfy the same inclusions. If
$U\subseteq W$, then
$T_a^*U\subseteq J_e(U)\subseteq J_e(W)$, so
$U\subseteq F_{a,e}(W)$. The induction over updates proves greatestness.
\end{proof}

\begin{corollary}[Relative completeness of discovery]
Any finite collection of degree-at-most-$d$ polynomials that vanishes initially
and admits preservation identities with multiplier degrees at most $e$ spans a
subspace of the computed result. In particular, the search does not miss a
certificate packet in this fixed fragment merely because its generators and
multipliers would have appeared bilinearly in a direct ansatz.
\end{corollary}
\begin{proof}
The span of the supplied collection is initially zero. Every linear combination
of its pullbacks is a bounded-degree polynomial combination of the same
collection, so its span satisfies the hypothesis of
Theorem~\ref{thm:ideal-greatest}.
\end{proof}

The corollary is about \emph{packets closed within the selected template and
multiplier bound}. It is not a decision procedure for all invariants, all
polynomial implications, or all true goals. A valid invariant may require
auxiliary generators above degree $d$, or preservation multipliers above degree
$e$, or a different certificate such as a radical-ideal argument.

\subsection{Worked nonlinear certificates}
For squaring on the diagonal, degree $d=1$, $e=0$ fails and $e=1$ succeeds:
\[
 p=y-x,\qquad p(x^2,y^2)=(x+y)p(x,y).
\]
More generally, for $T_k(x,y)=(x^k,y^k)$,
\[
 (y-x)\circ T_k=
 \left(\sum_{j=0}^{k-1}x^jy^{k-1-j}\right)(y-x).
\]
The archive synthesizes these packets for $k=2,\ldots,9$, using $e=k-1$.

An initial union is handled without choosing one component. Let initial states
be $(t,0)$ or $(0,t)$ and use the same squaring transition. In degree two,
$W_0=\Span\{xy\}$. Then
\[
 (xy)\circ T=x^2y^2=(xy)(xy).
\]
The degree-one multiplier limit fails, while $e=2$ recovers the packet. The
checker verifies both initial substitutions; checking only the first would
silently change the subject.

For the parabola $\sigma(t)=(t,t^2)$, take $p=y-x^2$. Then
\[
 p\circ T=y^2-x^4=(y+x^2)(y-x^2).
\]
Again $d=2,e=2$ succeeds. These examples exercise changing zero sets, unions of
initial images, and nonlinear multipliers without numerical sampling.

\subsection{What the checker actually certifies}
The certificate contains the exact problem, generator polynomials, and one
polynomial matrix per transition. Its checker verifies arities, canonical
rational coefficients, degree limits, all initial substitutions, and every
identity in \eqref{eq:poly-packet}. It calls neither nullspace computation nor
invariant synthesis. The replay driver replaces those functions with routines
that raise an exception if called.

An accepted certificate proves only the preservation claim represented by the
packet. It does \emph{not} certify that the packet is greatest. In particular,
an empty packet is a vacuously valid preservation certificate, not independently
checked evidence that no useful invariant exists. Certifying greatestness would
require additional row-space or elimination evidence for every restriction
step. Forge should not upgrade an empty returned packet to a mathematical
refutation of its goal.

\subsection{Cost, refinement, and a concrete production route}
There are $N=\binom{n+d}{d}$ template monomials and
$K=\binom{n+e}{e}$ multiplier monomials. At a stage with $k$ generators, the
restriction matrix has $k(1+K)$ columns. If the largest transition degree is
$t$, its rows range over monomials of degree at most $\max(dt,d+e)$, though
sparsity can make the actual support much smaller. At most $N$ strict
restrictions occur, and a simple pass scheduler uses at most $r(N+1)$ edge
checks. None of these bounds controls coefficient bit growth; rational
elimination and substitution expansion can dominate memory.

A useful initial scheduling order is $(d,e)=(1,0),(1,1),(2,0),(2,1),(2,2)$,
followed by problem-specific templates. This order is a proposed heuristic,
not a measured optimum. Increasing $e$ can recover certificates that a smaller
bound cannot express. However, a larger-bound run must begin with the full new
$W_0$, not the old final subspace: using the old answer as the only starting
space would prevent lost generators from reappearing.

Use sparse coefficient collection in the Lean-facing front end and an exact
rational matrix backend for the proposer. \code{python-flint}'s
\code{fmpq_mat} offers rational row reduction and solving as a concrete external
option; the included implementation uses \code{fractions.Fraction} to remain
self-contained~\cite{flint-matrix}. In Lean, reify into a small polynomial syntax,
prove its denotation correct, check coefficient identities, and instantiate a
generic reachability-invariant theorem. Mathlib's multivariate-polynomial
evaluation homomorphisms provide an appropriate semantic target~\cite{mvpoly}.
Neither external library is needed to trust a returned packet once the Lean
identity checker exists.

\subsection{Extension to control-flow locations}
\label{sec:locations}
A useful unimplemented next step is a family of generator spaces $W_\ell$, one
per control-flow location. Initial conditions restrict entry locations; a
location without an initial state starts with its full template. For an edge
$\ell\xrightarrow{T}k$, require
\[
 T^*W_k\subseteq J_e(W_\ell).
\]
The corresponding restriction is on the \emph{target} generator space:
\[
 W_k\leftarrow W_k\cap(T^*)^{-1}\bigl(J_e(W_\ell)\bigr).
\]
This remains monotone on the product of the subspace lattices and terminates
after at most the sum of their template dimensions in strict decreases. A
worklist rechecks an edge when either its source or target space changes.
The greatestness proof is the same containment induction as above.

An equality guard $g=0$ can contribute bounded multiples of $g$ to the
right-hand side, giving identities of the form
$p_k\circ T=\sum h_jp_{\ell j}+qg$. Inequality guards do not vanish and cannot
be inserted into this equation as if they did. Ignoring such guards gives a
sound over-approximation but may lose invariants; exploiting them requires a
separate proved implication. This control-flow extension is specified here,
not implemented in the single-location prototype.
